\subsection*{Motivation}

Chapter 9 gave us random variables; Chapter 10 gave us expectation and Jensen's inequality. We now ask: how much \emph{information} does an outcome carry, and how \emph{different} are two distributions on the same space? The answers --- entropy, cross-entropy, and the Kullback--Leibler divergence --- are not optional flavor. They \emph{are} the loss surface of every modern language model. Cross-entropy is the next-token training objective (Chapters 17, 25); KL is the trust-region regularizer in RLHF / DPO (Chapter 28); the entropy of the softmax is what people mean by ``sampling diversity.'' Every result below follows from one fact already proved: $-\ln$ is convex.

We use the natural log $\ln$ throughout for clean derivatives. Replacing $\ln$ by $\log_2$ multiplies every formula by $1/\ln 2$ and converts \emph{nats} to \emph{bits}; nothing else changes.

\subsection*{Definitions}

\begin{definition}[Self-information]
For a discrete random variable $X$ with PMF $p$ and $x \in \mathrm{supp}(p)$, the \emph{self-information} of the outcome $x$ is $I(x) := -\ln p(x)$. Rare outcomes carry more information; $I(x) \to \infty$ as $p(x) \to 0$.
\end{definition}

\begin{definition}[Shannon entropy]
The \emph{entropy} of $p$ is the expected self-information,
\[
H(p) := \mathbb{E}_{x \sim p}[I(x)] = -\sum_{x} p(x) \ln p(x),
\]
with the convention $0 \ln 0 = 0$, justified by $\lim_{t \downarrow 0} t \ln t = 0$.
\end{definition}

\begin{definition}[Joint and conditional entropy]
For a joint PMF $p_{X,Y}$,
\[
H(X, Y) = -\sum_{x,y} p_{X,Y}(x,y) \ln p_{X,Y}(x,y), \qquad H(Y \mid X) = -\sum_{x,y} p_{X,Y}(x,y) \ln p_{Y \mid X}(y \mid x).
\]
\end{definition}

\begin{definition}[Cross-entropy]
For PMFs $p, q$ on the same alphabet with $\mathrm{supp}(p) \subseteq \mathrm{supp}(q)$,
\[
H(p, q) := -\sum_{x} p(x) \ln q(x) = \mathbb{E}_{x \sim p}[-\ln q(x)].
\]
\end{definition}

\begin{definition}[Kullback--Leibler divergence]
\[
D_{\mathrm{KL}}(p \,\|\, q) := \sum_{x} p(x) \ln \frac{p(x)}{q(x)} = \mathbb{E}_{x \sim p}\!\left[\ln \frac{p(x)}{q(x)}\right].
\]
$D_{\mathrm{KL}}$ is \emph{not} symmetric and \emph{not} a metric.
\end{definition}

\begin{definition}[Mutual information]
For joint $p_{X,Y}$ with marginals $p_X, p_Y$,
\[
I(X; Y) := D_{\mathrm{KL}}(p_{X,Y} \,\|\, p_X \otimes p_Y) = \sum_{x,y} p_{X,Y}(x,y) \ln \frac{p_{X,Y}(x,y)}{p_X(x) p_Y(y)}.
\]
\end{definition}

\subsection*{Theorems and proofs}

\begin{theorem}[Gibbs' inequality / KL nonnegativity]
For PMFs $p, q$ on the same finite alphabet with $q(x) > 0$ wherever $p(x) > 0$,
\[
D_{\mathrm{KL}}(p \,\|\, q) \geq 0,
\]
with equality iff $p(x) = q(x)$ for every $x$ with $p(x) > 0$.
\end{theorem}

\begin{proof}
The function $\phi(t) = -\ln t$ is strictly convex on $(0, \infty)$ since $\phi''(t) = 1/t^2 > 0$. Compute, viewing the sum as $\mathbb{E}_{x \sim p}$:
\[
-D_{\mathrm{KL}}(p \,\|\, q) = \sum_{x : p(x) > 0} p(x) \ln \frac{q(x)}{p(x)} = -\,\mathbb{E}_{x \sim p}\!\left[\phi\!\left(\tfrac{q(x)}{p(x)}\right)\right].
\]
By Jensen's inequality (Chapter 10), $\mathbb{E}[\phi(Z)] \geq \phi(\mathbb{E}[Z])$ for convex $\phi$, hence
\[
-D_{\mathrm{KL}}(p \,\|\, q) \leq -\phi\!\left(\mathbb{E}_{x \sim p}\!\left[\tfrac{q(x)}{p(x)}\right]\right) = \ln\!\Bigl(\sum_{x : p(x) > 0} q(x)\Bigr) \leq \ln 1 = 0,
\]
using $\sum_x q(x) = 1$. Therefore $D_{\mathrm{KL}}(p \,\|\, q) \geq 0$. By strict convexity, Jensen is an equality iff $q(x)/p(x)$ is $p$-a.s.\ constant; since both $p$ and $q$ sum to $1$ on the support of $p$, that constant must be $1$, i.e.\ $p = q$ on $\mathrm{supp}(p)$.
\end{proof}

\begin{theorem}[Cross-entropy decomposition]
$H(p, q) = H(p) + D_{\mathrm{KL}}(p \,\|\, q).$
\end{theorem}

\begin{proof}
Direct algebra:
\[
H(p, q) = -\sum_x p(x) \ln q(x) = -\sum_x p(x) \bigl[\ln p(x) + \ln \tfrac{q(x)}{p(x)}\bigr] = H(p) + D_{\mathrm{KL}}(p \,\|\, q). \qedhere
\]
\end{proof}

\begin{corollary}
$H(p, q) \geq H(p)$, with equality iff $q = p$. In particular, $\min_q H(p, q) = H(p)$, attained at $q = p$.
\end{corollary}

\begin{theorem}[Maximum entropy on a finite alphabet]
For any PMF $p$ supported on a set of size $K$,
\[
H(p) \leq \ln K,
\]
with equality iff $p$ is the uniform distribution $u(x) = 1/K$.
\end{theorem}

\begin{proof}
Let $u$ denote the uniform distribution on the $K$-element support. Compute
\[
D_{\mathrm{KL}}(p \,\|\, u) = \sum_x p(x) \ln \frac{p(x)}{1/K} = -H(p) + \ln K \cdot \sum_x p(x) = \ln K - H(p).
\]
By Gibbs' inequality this is $\geq 0$, so $H(p) \leq \ln K$, with equality iff $p = u$.
\end{proof}

\begin{theorem}[Chain rule for entropy]
$H(X, Y) = H(X) + H(Y \mid X).$
\end{theorem}

\begin{proof}
Apply $\ln p_{X,Y}(x, y) = \ln p_X(x) + \ln p_{Y \mid X}(y \mid x)$:
\[
H(X, Y) = -\sum_{x,y} p_{X,Y}(x,y) \ln p_X(x) \;-\; \sum_{x,y} p_{X,Y}(x,y) \ln p_{Y \mid X}(y \mid x).
\]
Marginalizing $y$ in the first sum gives $-\sum_x p_X(x) \ln p_X(x) = H(X)$; the second sum is $H(Y \mid X)$ by definition.
\end{proof}

\begin{theorem}[Mutual information nonnegativity]
$I(X; Y) \geq 0$, with equality iff $X$ and $Y$ are independent.
\end{theorem}

\begin{proof}
By definition $I(X; Y) = D_{\mathrm{KL}}(p_{X,Y} \,\|\, p_X \otimes p_Y)$. Apply Gibbs (Theorem 1): this is $\geq 0$ with equality iff $p_{X,Y}(x, y) = p_X(x) p_Y(y)$ for all $x, y$ --- exactly independence.
\end{proof}

\begin{remark}
Combining the chain rule with the definition of $I$ yields the standard symmetric identity
\[
I(X; Y) = H(X) + H(Y) - H(X, Y) = H(Y) - H(Y \mid X) = H(X) - H(X \mid Y).
\]
Conditioning never increases entropy: $H(Y \mid X) \leq H(Y)$.
\end{remark}

\begin{definition}[Total variation distance]
\label{def:tv}
For two probability measures $P, Q$ on a measurable space $(\Omega, \mathcal{F})$,
\[
D_{TV}(P, Q) := \sup_{A \in \mathcal{F}} |P(A) - Q(A)|.
\]
For discrete distributions on a countable set $\mathcal{X}$, the supremum is attained on $A = \{x : P(x) \ge Q(x)\}$, giving the equivalent form
\[
D_{TV}(P, Q) = \tfrac{1}{2} \sum_{x} |P(x) - Q(x)| = \tfrac{1}{2}\|P - Q\|_1.
\]
\end{definition}

\begin{theorem}[Pinsker's inequality]
\label{thm:pinsker}
For any two probability measures $P, Q$ on a common space with $Q$ absolutely continuous with respect to $P$ (so $D_{KL}(P \| Q)$ is finite),
\[
D_{TV}(P, Q) \;\le\; \sqrt{\tfrac{1}{2}\, D_{KL}(P \| Q)}.
\]
\end{theorem}

\begin{proof}
\textbf{Step 1 (Bernoulli case).} Define, for $p, q \in (0,1)$, the binary KL and TV
\[
  d(p \| q) := p \ln\tfrac{p}{q} + (1-p)\ln\tfrac{1-p}{1-q}, \qquad d_{TV}(p, q) := |p - q|.
\]
We claim $d(p \| q) \ge 2 (p-q)^2$ for all $p, q \in (0,1)$. Fix $q$ and let $\varphi(p) := d(p \| q) - 2(p - q)^2$. Then $\varphi(q) = 0$, $\varphi'(p) = \ln\tfrac{p}{q} - \ln\tfrac{1-p}{1-q} - 4(p-q)$, and
\[
  \varphi''(p) = \tfrac{1}{p(1-p)} - 4 \;\ge\; 4 - 4 = 0,
\]
since $p(1-p) \le 1/4$. So $\varphi$ is convex in $p$ with minimum at $p = q$ (because $\varphi'(q) = 0$ and $\varphi$ is convex), giving $\varphi(p) \ge 0$ for all $p$. Hence $d(p \| q) \ge 2 (p - q)^2$.

\textbf{Step 2 (reduction to binary).} For any event $A \in \mathcal{F}$, set $p := P(A)$, $q := Q(A)$. Define a binary partition $\{A, A^c\}$ and consider the data-processing reduction: the map $\omega \mapsto \mathbf{1}\{\omega \in A\}$ pushes $P, Q$ to Bernoulli$(p)$ and Bernoulli$(q)$ respectively. By the data-processing inequality for KL divergence (which holds for any deterministic measurable map and follows directly from the log-sum inequality --- itself a consequence of Jensen applied to $\ln$),
\[
  D_{KL}(P \| Q) \;\ge\; d(p \| q).
\]
Combining with Step 1, $D_{KL}(P \| Q) \ge 2(p - q)^2 = 2 |P(A) - Q(A)|^2$. Take the supremum over $A$:
\[
  D_{KL}(P \| Q) \;\ge\; 2 \, D_{TV}(P, Q)^2,
\]
and rearrange to obtain the claimed inequality.
\end{proof}

\begin{remark}
Pinsker's inequality is the tool that turns KL bounds (which information-theoretic objectives like RLHF's $D_{KL}(\pi \| \pi_{\mathrm{ref}})$ control) into TV bounds (which couple to performance differences). Chapter~31's TRPO surrogate-improvement inequality uses it to convert a per-state KL penalty into a bound on the difference between state-visitation distributions $d^\pi$ and $d^{\pi'}$.
\end{remark}

\subsection*{Code sketch}

The notebook builds two categorical distributions $p, q$ on a 5-token alphabet, implements \texttt{entropy}, \texttt{cross\_entropy}, and \texttt{kl} from the definitions above, and verifies $H(p, q) = H(p) + D_{\mathrm{KL}}(p \,\|\, q)$ to machine precision. It samples $50$ random Dirichlet pairs (seed $0$) and confirms $D_{\mathrm{KL}} \geq 0$ for every pair, with $D_{\mathrm{KL}}(p \,\|\, p) = 0$. For $K = 8$, it confirms $H(\text{uniform}) = \ln 8$ and that $H(p) < \ln 8$ for $100$ non-uniform draws. Finally it constructs a correlated joint on $\{0, 1\}^2$, computes $I(X; Y)$ via both formulas, verifies agreement, and shows $I = 0$ for the product distribution.

\subsection*{Connection to LLMs}

A causal language model defines $p_\theta(x_t \mid x_{<t})$ via softmax (Chapter 9). On a corpus drawn from $p_{\mathrm{data}}$, the training objective is
\[
\mathcal{L}(\theta) = \mathbb{E}_{x \sim p_{\mathrm{data}}}\!\left[-\ln p_\theta(x_t \mid x_{<t})\right] = H(p_{\mathrm{data}}, p_\theta) = H(p_{\mathrm{data}}) + D_{\mathrm{KL}}(p_{\mathrm{data}} \,\|\, p_\theta).
\]
Theorem 2 makes precise why minimizing cross-entropy \emph{is} minimizing KL to the data: the entropy term $H(p_{\mathrm{data}})$ does not depend on $\theta$. Theorem 4 sets the absolute floor: a uniform LM over a $50000$-token vocabulary has entropy $\ln 50000 \approx 10.82$ nats, which is the loss every trained model improves on (Chapter 17, 25). In RLHF and DPO (Chapter 28), the policy update is regularized by $D_{\mathrm{KL}}(\pi_\theta \,\|\, \pi_{\mathrm{ref}})$ to keep the fine-tuned policy close to the SFT initialization --- Gibbs' inequality is exactly what makes that penalty a meaningful divergence. Finally, the entropy of the sampler's softmax, controlled by temperature, is the ``diversity'' knob in decoding.
